\documentclass[12pt,letterpaper]{report}
\usepackage{graphicx}
\usepackage{scrextend}
\usepackage{vmargin}
\usepackage{graphicx}
\usepackage{multirow}
\usepackage[utf8]{inputenc}
\usepackage[spanish]{babel}
\usepackage{multicol}
\usepackage{enumerate}
\usepackage{float}
\usepackage{amsmath, amsthm, amssymb, amsfonts}
\usepackage[usenames]{color}
\parindent=0mm
\pagestyle{empty}
\definecolor{miorange}{rgb}{0.91, 0.43, 0.0}
\begin{document}
\setmargins{2.5cm}      
{1.5cm}                     
{2cm}  
{24cm}                    
{10pt}                          
{1cm}                          
{0pt}                             
{2cm}
\begin{flushright}
\today
\end{flushright}
\vspace{-1.75cm}
\section*{Tarea Clase 9}
\subsection*{Obtenga las soluciones de las siguientes ecuaciones:}
\begin{enumerate}
\item Se tiene un rectángulo cuya altura mide 2cm más que su base y cuyo perímetro es igual a 24cm. Calcular las dimensiones del rectángulo.
\item El perímetro de un terreno rectangular es de 350 m. Sabiendo que el largo del terreno es el triple de su ancho, ¿cuáles son las dimensiones de la parcela?
\item La medida de los tres lados de un triángulo son tres números consecutivos. Si el perimetro del triángulo es 12 cm ¿cuánto mide cada lado?
\item Calcula las dimensiones de un rectángulo cuyo perimetro mide 48 cm, sabiendo que de largo mide el triple que de ancho.
\item Calcula las dimensiones de un rectangulo cuyo largo es siete cm más grande que el ancho y sus área es 70.
\end{enumerate}
\end{document}